Nesta seção, serão apresentados os métodos \gls{TOPSIS} e sua extensão \textit{\textit{Fuzzy}}-\gls{TOPSIS}. Esses métodos são responsáveis por ranquear um conjunto de alternativas com base em múltiplos critérios simultaneamente. Como etapa intermediária para a validação da arquitetura proposta neste trabalho, o \gls{TOPSIS} clássico foi inicialmente integrado ao algoritmo \gls{SAC} para a resolução do problema de descarregamento de tarefas, permitindo a avaliação dos candidatos e a redução do espaço de observação. No entanto, o objetivo final da pesquisa consiste em incorporar o \textit{\textit{Fuzzy}}-\gls{TOPSIS} para mitigar as incertezas e a imprecisão das medições inerentes aos cenários veiculares de alta mobilidade, modelando a incerteza linguística nos critérios de avaliação.

% ------------------------------------------------------------
\section*{O Algoritmo TOPSIS}
\label{subsec:topsis-algoritmo}
% ------------------------------------------------------------

O \gls{TOPSIS} é um método \gls{MCDM} que ranqueia alternativas pela proximidade relativa à \gls{PIS} (melhor valor possível em cada critério) e pela distância à \gls{NIS} (pior valor possível). Uma alternativa é preferível se estiver simultaneamente próxima da \gls{PIS} e distante da \gls{NIS}. A Figura~\ref{fig:topsis_geometrico} ilustra esse conceito de forma geométrica em um espaço bidimensional. Nela e na formulação matemática do método, a letra $d$ representa a distância euclidiana entre os pontos, os sobrescritos sinalizam qual é a fronteira de referência sendo medida (onde $^+$ simboliza o cálculo em direção à PIS, e $^-$ em direção à NIS), e o subscrito numérico refere-se a qual candidato $A_i$ está sendo avaliado. Observa-se que a alternativa $A_1$, por exemplo, apresenta uma distância menor em direção à métrica ideal positiva ($d^+_1$) e uma distância maior para a solução anti-ideal ($d^-_1$), tornando-se preferível em comparação direta com a alternativa $A_3$.

\begin{figure}[htb] 
\centering 
\begin{tikzpicture}[
    >=Latex,
    font=\sffamily,
    scale=1.2,
    dot/.style={circle, inner sep=0pt, minimum size=6pt},
    pis/.style={dot, fill=green!60!black, draw=black, thick},
    nis/.style={dot, fill=red!70, draw=black, thick},
    alt/.style={dot, fill=blue!50, draw=black, thick},
    line/.style={dashed, thick}
]
    % Box/Grid boundary
    \draw[dotted, thick, gray] (1, 1) rectangle (5, 4);

    % Axes
    \draw[->, thick, gray!80!black] (-0.5, 0) -- (6.5, 0) node[right, text=black] {Critério 1};
    \draw[->, thick, gray!80!black] (0, -0.5) -- (0, 5) node[above, text=black] {Critério 2};
    
    % Ideal Solutions
    \node[nis] (NIS) at (1, 1) {};
    \node[below left=2pt, font=\bfseries] at (NIS) {NIS ($A^-$)};
    
    \node[pis] (PIS) at (5, 4) {};
    \node[above right=2pt, font=\bfseries] at (PIS) {PIS ($A^+$)};
    
    % Alternatives
    \node[alt] (A1) at (4.2, 2.5) {};
    \node[above=3pt] at (A1) {$A_1$};
    
    \node[alt] (A2) at (3.2, 2.5) {};
    \node[above right=2pt] at (A2) {$A_2$};
    
    \node[alt] (A3) at (2.0, 3.2) {};
    \node[above left=1pt] at (A3) {$A_3$};

    % Distances for A1
    \draw[line, green!60!black] (A1) -- node[pos=0.4, right=2pt, font=\small, text=black] {$d^+_1$} (PIS);
    \draw[line, red!60!black] (A1) -- node[pos=0.4, below right=1pt, font=\small, text=black] {$d^-_1$} (NIS);
    
    % Distances for A3
    \draw[line, green!60!black] (A3) -- node[pos=0.4, above=2pt, font=\small, text=black] {$d^+_3$} (PIS);
    \draw[line, red!60!black] (A3) -- node[pos=0.4, left=2pt, font=\small, text=black] {$d^-_3$} (NIS);

\end{tikzpicture} 
\caption{Representação geométrica bidimensional do ranqueamento TOPSIS baseada em critérios de benefício.} 
\label{fig:topsis_geometrico}
\Fonte{Elaborado pelo autor.} 
\end{figure}

\begin{algorithm}[htb]
\LinesNumbered
\caption{Processo de Avaliação Multi-critério TOPSIS}
\label{alg:topsis}
\SetKwInOut{Input}{Entrada}
\SetKwInOut{Output}{Saída}
\Input{Matriz de decisão $\mathbf{X} \in \mathbb{R}^{m \times n}$ com $m$ alternativas em $n$ critérios; vetor normatizado de pesos $\mathbf{w}$.}
\Output{Vetor de proximidade $\mathbf{C}^*$ e alternativas ranqueadas decorrentemente.}
\BlankLine
\tcp{Passos 1 e 2: Normalização Vetorial e Ponderação}
\For{$j = 1$ \KwTo $n$}{
    \For{$i = 1$ \KwTo $m$}{
        $r_{ij} \gets x_{ij} \left/ \sqrt{\sum_{k=1}^{m} x_{kj}^2} \right.$\label{algline:norm}\;
        $v_{ij} \gets w_j \cdot r_{ij}$\label{algline:pond}\;
    }
    \BlankLine
    \tcp{Passo 3: Dimensionamento das Fronteiras PIS ($A^+$) e NIS ($A^-$)}
    \eIf{critério $j$ é de benefício}{
        $A^+_j \gets \max_i v_{ij}$\label{algline:pisnis_start}\;
        $A^-_j \gets \min_i v_{ij}$\;
    }{
        $A^+_j \gets \min_i v_{ij}$ \tcp*{Critério de custo (minimizar)}
        $A^-_j \gets \max_i v_{ij}$\label{algline:pisnis_end}\;
    }
}
\BlankLine
\tcp{Passos 4 e 5: Distâncias Euclidianas e Coeficiente de Similaridade}
\For{$i = 1$ \KwTo $m$}{
    $d^+_i \gets \sqrt{\sum_{j=1}^{n} (v_{ij} - A^+_j)^2}$\label{algline:distpis}\;
    $d^-_i \gets \sqrt{\sum_{j=1}^{n} (v_{ij} - A^-_j)^2}$\label{algline:distnis}\;
    $C^*_i \gets d^-_i / (d^+_i + d^-_i)$\label{algline:coeff}\;
}
\Return Ordenação decrescente de $\mathbf{C}^*$\label{algline:rank}\;
\end{algorithm}

De forma geral, o \gls{TOPSIS} organiza os dados em formato de matriz ($\mathbb{R}^{m \times n}$) e ranqueia $m$ alternativas candidatas simultaneamente considerando $n$ critérios através de um conjunto de passos. Tal procedimento, está abstraído no Algoritmo~\ref{alg:topsis}.
Ele funciona em três fases principais. Inicialmente, os dados de entrada são normalizados por coluna (linha~\ref{algline:norm}), o que unifica as métricas para que possuam a mesma escala, permitindo a comparação de critérios com naturezas e unidades diferentes. Em seguida, os valores normalizados são multiplicados por um vetor de pesos $\mathbf{w}$ que define a prioridade de cada critério (linha~\ref{algline:pond}).

Com os dados padronizados, o método identifica os pontos extremos. Para cada critério de benefício, busca-se o valor máximo; para critérios de custo, o valor mínimo. O agrupamento dos melhores valores forma a PIS, enquanto os piores formam a NIS (linhas~\ref{algline:pisnis_start}--\ref{algline:pisnis_end}).

Por último, o \gls{TOPSIS} calcula a distância Euclidiana de cada candidato até essas duas soluções ideais (linhas~\ref{algline:distpis} e~\ref{algline:distnis}). O coeficiente de proximidade $C^*_i \in [0,1]$ é obtido calculando a razão entre a distância à NIS e a distância total para ambos os extremos (linha~\ref{algline:coeff}). As alternativas são então ranqueadas ordenando $C^*_i$ decrescentemente (linha~\ref{algline:rank}). O uso específico da distância Euclidiana garante que alternativas com falhas simultâneas em vários critérios sofram penalidades maiores, valorizando soluções balanceadas. 

% Neste trabalho, conforme detalhado na Seção~\ref{subsec:topsis-motivacao}, a modelagem adota $n=4$ critérios principais para a tomada de decisão: latência ($\tau$), carga de \gls{CPU}, disponibilidade de banda ($C$) e consumo de energia ($E$).

% ------------------------------------------------------------
\section*{Extensão Fuzzy do TOPSIS}
\label{subsec:fuzzy-topsis}
% ------------------------------------------------------------

Em ambientes \gls{VEC}, por exemplo, as estimativas de qualidade do canal de comunicação, podem estar desatualizadas. Isto significa que o agente observa as condições do ambiente e ao tomar a decisão,  está tomando ações com base em informações incertas. Esta incerteza de medição abre espaço para uma generalização natural do \gls{TOPSIS}, a extensão \textit{\textit{fuzzy}}-\gls{TOPSIS}.

Nela, cada valor $x_{ij}$ da matriz de decisão é modelado como um TFN $\tilde{x}_{ij} = (L_{ij},\;M_{ij},\;U_{ij})$, onde $L_{ij}$ é o limite inferior, $M_{ij}$ é o valor modal e $U_{ij}$ é o limite superior do intervalo de confiança. A fuzzificação é realizada a partir da estatística da coluna $j$:

\begin{align}
    L_{ij} &= x_{ij} - \sigma_{\text{frac}} \cdot \hat{\sigma}_j, \label{eq:fuzz-lower} \\
    M_{ij} &= x_{ij}, \label{eq:fuzz-modal} \\
    U_{ij} &= x_{ij} + \sigma_{\text{frac}} \cdot \hat{\sigma}_j, \label{eq:fuzz-upper}
\end{align}

\noindent onde $\hat{\sigma}_j$ é o desvio padrão empírico da coluna $j$ e $\sigma_{\text{frac}} \in [0, 1]$ é um hiperparâmetro de fuzzificação. A defuzzificação pelo centroide recupera um valor escalar:

\begin{equation}
    \bar{x}_{ij} = \frac{L_{ij} + M_{ij} + U_{ij}}{3} = x_{ij},
    \label{eq:defuzz-centroid}
\end{equation}

\noindent que neste caso resulta exatamente no valor original (pois o TFN é simétrico em torno de $M_{ij}$), porém a estrutura \textit{\textit{fuzzy}} influencia os cálculos intermediários de PIS e NIS ao considerar os extremos do intervalo.

A justificativa para esta extensão é de natureza teórica: quando dois candidatos possuem escores $C^*$ muito próximos, qualquer perturbação de medição pode inverter o ranqueamento, gerando inconsistência na política do agente~\cite{farimani2024}. A extensão \textit{\textit{fuzzy}} modela este intervalo de incerteza via TFN, prevenindo inversões instáveis de ranqueamento ao considerar explicitamente os extremos do intervalo de confiança de cada medição.

A propriedade de unificação paramétrica é central para a arquitetura proposta: quando $\sigma_{\text{frac}} = 0$, tem-se $L_{ij} = M_{ij} = U_{ij} = x_{ij}$, e o método reduz-se exatamente ao \gls{TOPSIS} padrão. Esta parametrização permite que o mesmo módulo opere tanto no modo determinístico quanto no modo \textit{\textit{fuzzy}}, com a transição controlada por um único hiperparâmetro — sem qualquer alteração estrutural da arquitetura. Na presente proposta, adota-se $\sigma_{\text{frac}} = 0{,}0$, correspondendo ao \gls{TOPSIS} determinístico padrão. A extensão \textit{\textit{fuzzy}} está disponível como grau de liberdade de projeto para adaptação a ambientes com maior incerteza de medição, sem custo de redesenho do sistema.
